% gz-p10-06-1: h(x) = ln(x)/x 图象 + 极大值 h(e) = 1/e + 三条水平线 y=k
\documentclass[tikz,border=4pt]{standalone}
\usepackage{ctex}
\usepackage{amsmath}
\usepackage{pgfplots}
\pgfplotsset{compat=1.18}

\begin{document}
\begin{tikzpicture}
\begin{axis}[
  width=12cm, height=7cm,
  axis lines=center,
  xlabel={$x$}, ylabel={$y$},
  every axis x label/.style={at={(axis cs:10.2,0)}, anchor=west},
  every axis y label/.style={at={(axis cs:0,0.55)}, anchor=south},
  xmin=0, xmax=10.2,
  ymin=-0.9, ymax=0.55,
  xtick={2.718},
  xticklabels={$e$},
  ytick={0.368},
  yticklabels={$\frac{1}{e}$},
  tick style={color=black},
  ticklabel style={font=\small},
  clip=true,
]
  % 曲线 y = ln(x) / x，x > 0
  \addplot[black, thick, domain=0.25:10.0, samples=200] {ln(x)/x};
  \node[right, black] at (axis cs:6.5, 0.27) {$y=\frac{\ln x}{x}$};

  % 极大值 (e, 1/e)
  \addplot[only marks, mark=*, mark size=2.5pt, red] coordinates {(2.718, 0.368)};
  \node[above, red, font=\small] at (axis cs:2.718, 0.368) {极大$(e,\frac{1}{e})$};

  % 三条水平线 y = k
  % k = 1/e (=0.368): 1 个交点（切于极大值）
  \addplot[blue, dashed, thick, domain=0.25:10.0] {0.368};
  \node[left, blue, font=\footnotesize] at (axis cs:0.5, 0.41) {$k=\frac{1}{e}$: 1 个交点};

  % k = 0.25: 2 个交点
  \addplot[red, dashed, thick, domain=0.25:10.0] {0.25};
  \node[right, red, font=\footnotesize] at (axis cs:8.0, 0.28) {$k=0.25$: 2 个交点};

  % k = -0.2: 1 个交点（左侧）
  \addplot[purple, dashed, thick, domain=0.25:10.0] {-0.2};
  \node[right, purple, font=\footnotesize] at (axis cs:5.5, -0.25) {$k=-0.2$: 1 个交点};

  % 虚线
  \draw[gray, dashed] (axis cs:2.718, 0.368) -- (axis cs:2.718, 0);
  \draw[gray, dashed] (axis cs:2.718, 0.368) -- (axis cs:0, 0.368);
\end{axis}
\end{tikzpicture}
\end{document}
